\documentclass[11pt]{article}

\usepackage[margin=1in]{geometry}
\usepackage{amsmath,amssymb,mathtools}
\usepackage{braket}
\usepackage[colorlinks=true,linkcolor=blue]{hyperref}
\newcommand{\lb}{\left(}
\newcommand{\rb}{\right)}
\newcommand{\Tr}{\mathrm{tr}}
\newcommand{\mA}{\mathcal{A}}
\newcommand{\mE}{\mathcal{E}}
\newcommand{\mH}{\mathcal{H}}
\newcommand{\mL}{\mathcal{L}}
\newcommand{\nn}{\nonumber}

\begin{document}


\section{Mathematical structure in the group case}
\label{sec:group-case-appendix}

In a general locally compact group $G$, we have a left-invariant Haar measure $d\mu_L(g)$ which allows us to study the group $G$ through the lens of measurable functions that live on it. The Haar measure allows us to define $L_p$-norms
\begin{align}
    \|f\|_{p,L}=\lb\int d\mu_L(g) |f(g)|^p\rb^{1/p}
\end{align}
We also have a right-invariant Haar measure $d\mu_R(g)$ we satisfies
\begin{align}
\mu_L(g)=\mu_R(g^{-1})=\Delta(g) \mu_R(g)    
\end{align}
where $\Delta(g)$ is called the modular function of $G$. Since $\|f\|_{p,L}=\|f\|_{p,R}$ we simplify our notation by defining $\|f\|_p\equiv \|f\|_{p,L}$ and we will only work with the left-invariant Haar measure $\mu(g)\equiv \mu_L(g)$. To keep track of the right-Haar measure, we introduce the anti-pode map $S:g\to g^{-1}$ which is self-adjoint and satisfies $S=S^\dagger=S^{-1}$. 

For simplicity, we will focus on compact groups, which ensures that the left and right Haar measures match and $\Delta(g)=1$ (the group is unimodular), and the Haar measure is finite. The $L^\infty(G)$ is the isometric dual of $L^1(G)$ via the Hölder pairing,
\begin{equation}
    \langle f, h \rangle \equiv \int_G f(x) h(x) \, d\mu(x)\ .
\end{equation}
For compact groups, we have the inclusion of closed vector spaces
\begin{align}\label{hierLp}
    L^\infty(G)\subseteq L^2(G)\subseteq L^1(G)
\end{align}
Note that $L^2(G)$ is self-dual, whereas $L^\infty(G)$ is the Holder dual of  $L^1(G)$ \footnote{In this note, we will be primarily concerned with finite groups where this hierarchy simplifies to $\mathbb{C}(G)\simeq L^1(G)\simeq L^2(G)\simeq L^\infty(G)$.}. 
% To make Holder duality manifest, we choose the notation 
% \begin{align}
%     &L^\infty(G)\ni \ket{f}=\sum_g f(g) \ket{g}\nn\\
%     &L^1(G)\ni \bra{h}=\sum_g h(g)\bra{g}\ .
% \end{align}
To turn these linear vector spaces into algebras, there are two choices of multiplication we can make:
\begin{align}
    &(f_1\cdot f_2)(g)=f_1(g)f_2(g)\nn\\
    &(f_1*f_2)(g)=\int d\mu(h) \: f_1(h)f_2(h^{-1}g)
\end{align}
These multiplications result in two different algebras: the Abelian algebra of $L^\infty(G)$, which encodes the geometry of $G$, and the convolution algebra $L^1(G)$, which is non-Abelian if and only if $G$ is non-Abelian. 

\paragraph{Algebra of observables $L^\infty(G)$:} The first multiplication is manifestly commutative, and since $\|f_1\cdot f_2 \|_\infty\leq \|f_1\|_\infty\|f_2\|_\infty$, the vector space $L^\infty(G)$ is closed with respect to this multiplication, and hence an Abelian algebra. This Abelian algebra encodes the {\it geometry of group $G$} \footnote{For any measurable space $X$, we have an Abelian algebra $L^\infty(X)$, and a classical result establishes that to every Abelian algebra one can associate a measurable space. We will refer to this as classical Gelfand duality.}. In an analogy with classical mechanics that interprets $G$ as the phase space, functions $L^\infty(G)$ have the interpretation of classical observables. 

\paragraph{Algebra of states $L^1(G)$ (pre-dual):} In contrast, $L^1(G)$ does not form an algebra under pointwise multiplication due to non-integrable singularities in the squared functions. Instead, it forms a Banach algebra under left convolution. By Young's convolution inequality, $\|f * g\|_1 \leq \|f\|_1 \|g\|_1$, ensuring closure. In the left regular representation, these $L^1$ functions define bounded integral operators that smear the group action over the manifold.

There are a few comments:

\begin{itemize}
    \item {\bf Dual group $\hat{G}$:} For Abelian groups, $L^1(G)$ is also Abelian and encodes the geometry of the Fourier space to $G$ that we denote by $\hat{G}$. In the non-Abelian case, we will say that $L^1(G)$ encodes a {\it dual group} $\hat{G}$ which is neither a measurable space nor a group. It is sometimes called a {\it quantum group} (as an example of {\it quantum geometry}). By non-Abelian Gelfand duality, the irreducible representations $r$ of $G$ play the role of {\it quantum points} in the {\it quantum geometry}.

\item {\bf Right convolution algebra:} We can also define a right-convolution multiplication: $(f_1*_R f_2)=f_2* f_1$. With this multiplication, we obtain a different Banach space $L_R^1(G)$, and we can show explicitly that all operators in $L^1(G)$ and $L^1_R(G)$ commute. In fact, they are isometric with the isomorphism established by the anti-pode map $S:g\to g^{-1}$ that swaps the left-convolution with the right-convolution.

\end{itemize}

The algebras above can be represented as subalgebras of operators in the regular representation $L^2(G)$:
\begin{align}
&\iota_1(L^1(G))\subset B(L^2(G))\nn\\
&\iota_\infty(L^\infty(G))\subset B(L^2(G))\ .
\end{align}
with embeddings (representations)
\begin{align}
    &\iota_\infty(f)=\sum_g f(g) e_g,\qquad e_g\equiv \ket{g}\bra{g}\nn\\
    &\iota_1(f)=\sum_g f(g) U_g, \qquad U_h=\sum_g \ket{h g}\bra{g}\ .
\end{align}
The right convolution algebra $L^1_R(G)\simeq L^1(G)$ is embedded by the isometry
\begin{align}
    \iota_1^R(f)=\sum_g f(g) U'_g, \qquad U'_h\equiv \sum_g \ket{g h^{-1}}\bra{g}=\sum_g \ket{g}\bra{g h}=S U_g S
\end{align}
where $S$ is the anti-pode. Note that $U_g$ and $U'_g$ are unitary, and we can write 
\begin{align}
    L^1(G)=\mathbb{C}\rtimes_{\iota_1} G
\end{align}
Note that $U_{g_1}U_{g_2}=U_{g_1g_2}$ implies that the multiplication of two functions $\delta(g-g_1)$ and $\delta(g-g_2)$ should result in $\delta(g_1g_2)$ which fixes the multiplication to be convoltuion:
\begin{align}
 \delta(g-g_1)*\delta(g-g_2)=\delta(g-g_1g_2).   
\end{align}
Since $S$ acts as a unitary on $B(L^2(G))$ we also have a second representation for $L^\infty(G)$. We inroduce the notation
\begin{align}
    &\iota^\pm_\infty(f)=\sum_g f(g) e_g^\pm, \qquad e_g^+=e_g,\qquad  e_g^-=\ket{g^{-1}}\bra{g^{-1}}\nn\\
    &\iota^\pm_1(f)=\sum_g g(h) U_g^\pm, \qquad U_g^+=U_g, \qquad U_g^-=U'_g\nn\\
    &\iota_1^\pm=S\circ \iota_1^\mp\circ S,\qquad \iota_\infty^\pm=S\circ \iota_\infty^\mp\circ S
\end{align}
and we have the relations
\begin{align}
    &U^\pm_g e_h^\pm=e^\pm_{g h}U_h^\pm,,\qquad U^\mp_g e_h^\pm=e^\pm_{h g^{-1}} U_g^{\mp}
\end{align}

As we saw $L^1_R(G)$ is in the relative commutant of the inclusion $\iota_1$. 
In fact, we have
\begin{align}
    &\iota_1(L^1(G))'\simeq Z(\iota_1(L^1(G)))\vee \iota_1^R(L^1_R(G))\nn\\
    &\iota_1(L^1(G))\vee \iota_1^R(L^1_R(G))\simeq B(L^2(G))\ .
\end{align}
Furthermore, we have
\begin{align}
    L^1(G)\vee L^\infty(G)=B(L^2(G))\ .
\end{align}

For a finite group, it is natural to use the GNS Hilbert space representation of $B(L^2(G))$ in the tracial state
\begin{align}
    &\omega(a)=\frac{1}{|G|}\sum_{g\in G}\braket{g| a |g}\nn\\
    &\pi_\omega(1)=\ket{1}\equiv\frac{1}{\sqrt{|G|}}\sum_{g\in G}\ket{g}\ket{g}\nn\\
    &\pi_\omega(a)=(a\otimes 1)\ket{1}\ .
\end{align}
Then, we define the projections $P_1^\pm$ and $P_\infty=P_\infty^\pm$
\begin{align}
    &\mH_1^\pm=P_1^\pm\mH\equiv \overline{\iota_1^\pm(f)\ket{1}}\nn\\
    &\mH_\infty^\pm=P_\infty^\pm\mH\equiv \overline{\iota_1^\pm(f)\ket{1}}\ .
\end{align}
These projections are the {\it canonical endormophism} for the inclusion $\iota_1(L^1(G))\subset B(L^2(G))$.

\begin{align}
    e_g\to \ket{g}\ket{g},\qquad U_g\to \sum_h \ket{g h}\ket{h}, \qquad S\ket{g}\ket{h}=\ket{h}\ket{g}\ .
\end{align}


Therefore, we obtain maps 
\begin{align}
    &\iota_\infty^\dagger\circ \iota_1:L^1(G)\to L^\infty(G)\nn\\
    &\iota_1^\dagger\circ \iota_\infty:L^\infty(G)\to L^1(G)\ .
\end{align}
For compact groups, due to the hierarchy eq (\ref{hierLp}),  $\iota_\infty^\dagger\circ \iota_1$ is an embedding ($*$-homomorphism), whereas for finite groups $G$, it is an isomorphism realized by a unitary map that does a non-Abelian Fourier transform, swapping point-by-point multiplication in $G$ by convolution in $G$. 

We have a natural action of $G$ on the algebra $L^1(G)$. The subalgebra of $L^1(G)$ invariant under this action the algebra of class functions. That is to say
\begin{align}
    U_h \lb \sum_g f_g U_g\rb U_{h^{-1}}=\sum_g f_g U_{h g h^{-1}}=\sum_g f_g U_g
\end{align}
if and only if $f_g=f_{h g h^{-1}})$. In this case, this subalgebra is isomorphic to $Z(L^1(G))$. 

We aim to study the following inclusion of algebras
\begin{align}
 Z(L^1(G))\subset L^1(G)\subset B(L^2(G))\ .
\end{align}
Note that here $Z(L^1(G))$ is the same as the 


There can be many embeddings of $L^2(G)$ inside $B(L^2(G))$. 


To keep our discussion generalizable beyond finite groups, we choose one embedding for $L^\infty(G)$

There is always a geometric embedding 
\begin{align}
    &\iota_\infty:\bra{g}\to e_g=\ket{g}\bra{g}\nn\\
    &\iota_\infty(\bra{f})=\sum_g f(g)\ket{g}\bra{g}
\end{align}
that results in point-by-point multiplication and an Abelian algebra. 

An alternative embedding is provided by
\begin{align}
    &\iota_1:\ket{g}\to T_g=\sum_h \ket{g h}\bra{h}\nn\\
    &\iota_1(\ket{f})=\sum_{gh}f(g)\ket{g h}\bra{h}
\end{align}
That results in left convolution multiplication and a non-Abelian algebra if and only if $G$ is non-Abelian.

It is convenient to choose the tracial state on $B(L^2(G))$ and construct the GNS Hilbert space by building the representation:


This embedding combined with its dual, gives the conditional expectation
\begin{align}
(\iota_\infty\circ\iota_\infty^\dagger)(a\in B(L^2(G)))=\sum_g \braket{g|a|g}=\textbf{Tr}_{L^2(G)}(a) 
\end{align}


\subsection{Groups and their duals}

We take $G$ to be a finite group. The analogous statements for compact groups
follow by replacing sums with Haar integrals. The regular representation acts on
$L^2(G)$, with orthonormal basis $\{\ket g\}_{g\in G}$ and inner product
\begin{eqnarray}
    \braket{f_1|f_2}=\sum_{g\in G}f_1^*(g) f_2(g).
\end{eqnarray}

The left and right regular actions are
\begin{eqnarray}
    U_h\ket{g}=\ket{h g}, \qquad U'_h\ket{g}=\ket{g h^{-1}}.
\end{eqnarray}
Equivalently,
\begin{eqnarray}
    U_h=\sum_{g\in G}\ket{h g}\bra{g},\qquad
    U'_h=\sum_{g\in G}\ket{g h^{-1}}\bra{g}
        =\sum_{g\in G}\ket{g}\bra{g h}\ .
\end{eqnarray}
The two actions commute:
\begin{eqnarray}
    U_{h_1}U'_{h_2}=\sum_{g\in G}\ket{h_1 g h_2^{-1}}\bra{g}=U'_{h_2}U_{h_1}\ .
\end{eqnarray}
The sets $\{U_h\}$ and $\{U'_h\}$ generate commuting representations of
$G$ on $L^2(G)$.
The von Neumann algebra generated by $\{U_h\}$ is 
\begin{align}
    &\mathcal{L}(G)=\{U_g:g\in G\}''\subset B(L^2(G))
\end{align}
and its commutant is 
\begin{align}
    &\mathcal{L}(G)'=\{U'_g:g\in G\}''\ .
\end{align}
This is the left/right representation duality. Note that the center of $\mathcal{L}(G)$ is already included inside $\mathcal{L}(G)'$:
\begin{align}
    Z(L^1(G))=\mathcal{L}(G)\cap \mathcal{L}(G)'
\end{align}
Note that for finite groups we have $\mathcal{L}(G)\simeq L^1(G)\simeq L^2(G)\simeq \mathbb{C}(G)$.
Applying the Haar average over $G$ projects $L^1(G)$ to $Z(L^1(G))$:
\begin{align}
    &\mathcal{E}(a)=\frac{1}{|G|}\sum_g U_g a U_g^\dagger\nn\\
    &\mathcal{E}:L^1(G)\to Z(L^1(G))\ .
\end{align}

By the Peter--Weyl decomposition,
\begin{align}
    L^2(G)=\oplus_{r\in\hat{G}} \mH_r\otimes \mH_r^*
\end{align}
where $\hat G$ denotes the set of irreducible representations of $G$ and
$d_r=\dim\mH_r$. The word ``dual group'' should be used with care:
$\hat G$ is itself a group only when $G$ is Abelian. In this decomposition,
\begin{align}
    &\mathcal{L}(G)=\oplus_{r\in\hat G} B(\mH_r)\otimes 1_{\mH_r^*}\nn\\
    &\mathcal{L}(G)'=\oplus_{r\in\hat G} 1_{\mH_r}\otimes B(\mH_r^*)\ .
\end{align}
Thus, the Wedderburn block of the left regular group algebra associated with
$r$ is precisely $B(\mH_r)\cong M_{d_r}(\mathbb C)$; the factor $\mH_r^*$ is
the multiplicity space with which this block appears in the regular
representation.
An element of the group algebra is $X_f=\sum_h f(h)U_h$. Multiplication in
$\mathcal L(G)$ is left convolution:
\begin{align}
    &X_{f_1}X_{f_2}=X_{f_1*f_2}\,,\nn\\
    &(f_1*f_2)(h)=\sum_{g\in G} f_1(h g^{-1})f_2(g)\ .
\end{align}
Similarly, multiplication in $\mathcal{L}(G)'$ is right convolution. The
conditional expectation onto the left group algebra is obtained by averaging over
the right regular action,
\begin{align}
    \mE_{\mathcal L}(T)=\frac{1}{|G|}\sum_{g\in G}U'_g T (U'_g)^\dagger,
    \qquad
    \mE_{\mathcal L}:B(L^2(G))\to\mathcal L(G).
\end{align}
Averaging over the left regular action instead projects onto the commutant
$\mathcal L(G)'$. The normalized trace on $B(L^2(G))$, restricted to
$\mathcal L(G)'$, gives the Plancherel state
\begin{align}
    \tau(X')=\frac{1}{|G|}\Tr_{L^2(G)}(X')
    =\sum_{r\in\hat G}\frac{d_r}{|G|}
      \Tr_{\mH_r^*}(X'_r),
    \qquad
    X'=\oplus_r\,1_{\mH_r}\otimes X'_r .
\end{align}
On the central projection associated with $r$, this state has weight
$d_r^2/|G|$.



\subsection{Fourier transform, GNS quantization, and fusion}

Using the Plancherel state on $\mathcal L(G)'$, the GNS Hilbert space is
canonically isomorphic to $L^2(G)$. With the convention
$\Lambda(U'_g)=\ket g$, the identity element of the right group algebra is
mapped to $\ket e$. In the representation basis this is
\begin{align}
    \Lambda(1')=\ket e
    =\sum_{r\in\hat G}\sum_{i=1}^{d_r}
      \sqrt{\frac{d_r}{|G|}}\ket{r;i,i}.
\end{align}
More generally, the non-Abelian Fourier transform is
\begin{align}
    \ket{r;i,j}
    =\sqrt{\frac{d_r}{|G|}}\sum_{g\in G}
      \overline{U^r_{ij}(g)}\,\ket g,
    \qquad
    \ket g
    =\sum_{r\in\hat G}\sum_{i,j=1}^{d_r}
      \sqrt{\frac{d_r}{|G|}}U^r_{ij}(g)\ket{r;i,j}.
\end{align}
Here $U^r(g)$ is a unitary matrix for the irrep $r$. Modular conjugation
exchanges the two regular algebras, $J\mathcal{L}(G)J=\mathcal{L}(G)'$. In the group-theoretic setup, this is the antipode that sends $g\to g^{-1}$.

The GNS map turns multiplication in the group algebra into a bilinear product on
the Hilbert space:
\begin{align}
    \ket{g_1}*\ket{g_2}=\ket{g_1g_2}.
\end{align}
This product is convolution. Under the Fourier transform it becomes blockwise
matrix multiplication:
\begin{align}
    \mathbb C[G]\simeq \bigoplus_{r\in\hat G}B(\mH_r),
    \qquad
    E^r_{ij}E^s_{k\ell}
    =\delta_{rs}\delta_{jk}E^r_{i\ell},
\end{align}
where $E^r_{ij}$ are matrix units in the $r$ block. Thus the convolution
algebra remembers the Wedderburn blocks of the regular representation, but by
itself it does not encode the tensor-product fusion rules.

The fusion rules of $\mathrm{Rep}(G)$ enter through the additional Hopf-algebra
structure, namely the coproduct
\begin{align}
    \Delta(U_g)=U_g\otimes U_g.
\end{align}
Equivalently, they appear in pointwise products of matrix coefficient functions.
The tensor product of irreducible representations decomposes as
\begin{align}
    \mH_r\otimes\mH_s\cong
    \bigoplus_{t\in\hat G} N_{rs}^t\,\mH_t,
\end{align}
with Clebsch--Gordan intertwiners
\begin{align}
    C^{t,\alpha}_{rs}:\mH_r\otimes\mH_s\to\mH_t,
    \qquad \alpha=1,\cdots,N_{rs}^t.
\end{align}
This is the point at which the categorical fusion data of $\mathrm{Rep}(G)$
enters.

\subsection{Coactions and crossed products}

The discussion above shows that the group algebra $\mathcal L(G)$ acts on
$L^2(G)$ by left convolution, while its commutant $\mathcal L(G)'$ acts by right
convolution. This left-right structure can be reformulated in a way that is more
flexible and admits generalizations, namely in terms of coactions and crossed
products.

The crossed product of the trivial algebra $\mathbb C$ by the group action is
the group algebra:
\begin{align}
    \mathcal{L}(G)\simeq \mathbb{C}\rtimes G.
\end{align}
Adjoining the algebra $C(G)$ of multiplication operators,
\begin{align}
    M_f\ket g=f(g)\ket g,\qquad f\in C(G),
\end{align}
to the left convolution algebra generates the full matrix algebra:
\begin{align}
    \mathcal L(G)\vee C(G)=B(L^2(G)).
\end{align}
For Abelian $G$, $C(G)\simeq\mathcal L(\hat G)$, and this statement is the
finite-dimensional form of crossing again by the dual group. The dual action is
most simply written as
\begin{align}
    \alpha_\chi(U_g)=\chi(g)U_g,
    \qquad \chi\in\hat G .
\end{align}
For non-Abelian $G$, $\hat G$ is not an ordinary group. The replacement is the
Hopf-algebra coproduct, or coaction,
\begin{align}
    \delta:\mathcal L(G)\to\mathcal L(G)\otimes\mathcal L(G),
    \qquad
    \delta(U_g)=U_g\otimes U_g.
\end{align}
On the Hilbert space one can also use the corresponding comodule map
\begin{align}
    V:L^2(G)\to L^2(G)\otimes L^2(G),
    \qquad
    V\ket g=\ket g\otimes\ket g,
\end{align}
which is the vector-space version of the same group-like structure.

\subsection[Example of Z2]{Example of $\mathbb{Z}_2$}

When the underlying physical space is restricted to the simplest possible finite group, $G = \mathbb{Z}_2 = \{0, 1\}$, the abstract machinery of operator algebras collapses into the foundational mathematics of quantum computing. The function spaces become finite-dimensional sequences, and the duality between classical geometry and quantum symmetry manifests exactly as the relationship between the Pauli $Z$ and Pauli $X$ operators.

For $G = \mathbb{Z}_2$, the physical state space $L^2(\mathbb{Z}_2)$ is simply the 2-dimensional Hilbert space $\mathbb{C}^2$. Any function on this group is a sequence of two complex numbers: $f = (f(0), f(1))$. We define the standard coordinate basis vectors:
\begin{itemize}
    \item $\delta_0 = (1, 0)$ corresponding to the state $|0\rangle$.
    \item $\delta_1 = (0, 1)$ corresponding to the state $|1\rangle$.
\end{itemize}

\paragraph{$L^\infty(\mathbb{Z}_2)$ (Pointwise Multiplication)/Classical geometry:} If we treat $\mathbb{Z}_2$ purely as a static classical space, we equip $\mathbb{C}^2$ with pointwise multiplication: $(f \cdot h)(g) = f(g)h(g)$.
Under this product, the basis vectors act as mutually orthogonal projections: $\delta_0 \cdot \delta_1 = 0$. Represented as operators acting on $\mathbb{C}^2$, they form the algebra of diagonal matrices:
\begin{equation}
    \delta_0 \mapsto \begin{pmatrix} 1 & 0 \\ 0 & 0 \end{pmatrix}, \quad \delta_1 \mapsto \begin{pmatrix} 0 & 0 \\ 0 & 1 \end{pmatrix}
\end{equation}
This commutative algebra is $L^\infty(\mathbb{Z}_2)$. It represents classical coordinate measurements and is generated by the \textbf{Pauli $Z$} operator, which corresponds to the function $f(g) = (-1)^g = \delta_0 - \delta_1$.

\paragraph{$L^1(\mathbb{Z}_2)$, (Convolution multiplication)/Quantum Symmetry:} If we treat $\mathbb{Z}_2$ as a group of dynamical translations, we equip $\mathbb{C}^2$ with discrete convolution: $(f * h)(x) = \sum_y f(y)h(y^{-1}x)$. 
Because the group operation is addition modulo 2, the basis vectors mix rather than annihilate: $\delta_1 * \delta_1 = \delta_0$. Represented as operators acting via left convolution, they form the algebra of permutation matrices:
\begin{equation}
    \delta_0 \mapsto \begin{pmatrix} 1 & 0 \\ 0 & 1 \end{pmatrix} = I, \quad \delta_1 \mapsto \begin{pmatrix} 0 & 1 \\ 1 & 0 \end{pmatrix} = X
\end{equation}
This commutative algebra is the group algebra $\mathcal{L}(\mathbb{Z}_2)$. It represents active spatial translations and is exactly the algebra generated by the \textbf{Pauli $X$} operator.

\paragraph{Duality:} The generalized Fourier transform, which bridges the pointwise geometry of $L^\infty(\mathbb{Z}_2)$ and the convolutional symmetry of $\mathcal{L}(\mathbb{Z}_2)$, is physically realized as the \textbf{Hadamard gate} ($H$). It acts as the perfect dictionary between the two dual algebras:
\begin{equation}
    H Z H = X, \quad H X H = Z
\end{equation}
Furthermore, the fundamental theorem of crossed products---which states that the geometry and symmetries jointly generate the full universe of bounded operators on the phase space---reduces to the standard fact that the $Z$ algebra and $X$ algebra jointly span the space of all $2 \times 2$ matrices:
\begin{equation}
    \mathcal{L}(\mathbb{Z}_2) \vee L^\infty(\mathbb{Z}_2) = \mathcal{B}(\mathbb{C}^2) \cong M_2(\mathbb{C})
\end{equation}

For $G=\mathbb{Z}_2=\{e,g\}$, we have $L^2(G)=\mathbb C^2$ spanned by
$\ket e$ and $\ket g$. The nontrivial left and right regular actions coincide:
$U_g=U'_g=X$. Thus $\mathcal L(G)=\mathcal L(G)'$ is the Abelian algebra
generated by $\{\mathbb I,X\}$:
\begin{align}
    \mathcal L(\mathbb Z_2)
    =\{c_+\ket{+}\bra{+}+c_-\ket{-}\bra{-}\},
    \qquad
    \ket{\pm}=\frac{1}{\sqrt 2}(\ket e\pm\ket g).
\end{align}
The Plancherel measure is uniform on the two irreducible sectors. In this
Abelian example, convolution diagonalizes in the Fourier basis,
\begin{align}
    P_\pm=\frac{1}{2}(\mathbb I\pm X),\qquad
    P_\pm P_\pm=P_\pm,\qquad P_+P_-=0.
\end{align}
The tensor-product fusion rule of irreducible representations is a separate
statement:
\begin{align}
    (+)\otimes(-)=(-),\qquad
    (-)\otimes(-)=(+).
\end{align}

\subsection{Operator Algebraic Structures on Locally Compact Groups}

The discussion of finite groups above naturally embeds into a broader framework. For a locally compact group $G$, the construction of classical and quantum observable algebras begins with the ``seed'' space of continuous complex-valued functions with compact support, $C_c(G)$. By equipping this space with different multiplication rules and completing it under their corresponding natural norms, we generate two distinct mathematical branches: the classical geometry and the quantum symmetry.

To elevate these spaces to the level of exact quantum measurements (von Neumann algebras), they are represented as operators acting on the physical phase space $L^2(G)$, and closed under the weak operator topology (the double commutant).

\paragraph{1. The Classical Geometry Branch (Commutative)}
This branch treats $G$ strictly as a classical, static geometrical manifold.
\begin{itemize}
    \item \textbf{Multiplication:} Pointwise, $(f \cdot g)(x) = f(x)g(x)$.
    \item \textbf{Natural Norm:} The supremum norm, $\|f\|_\infty = \sup_{x \in G} |f(x)|$.
    \item \textbf{Norm Closure (C*-algebra):} Completing $C_c(G)$ under the supremum norm yields $C_0(G)$, the space of continuous functions vanishing at infinity.
    \item \textbf{Weak Closure (von Neumann Algebra):} Representing $C_0(G)$ on $L^2(G)$ via pointwise multiplication and taking the double commutant yields $L^\infty(G)$, the algebra of all essentially bounded measurable functions. This represents the macroscopic classical observables.
\end{itemize}

\paragraph{2. The Quantum Symmetry Branch (Non-Commutative)}
This branch treats $G$ as a dynamical set of active translations or symmetries.
\begin{itemize}
    \item \textbf{Multiplication:} Left convolution, $(f * g)(x) = \int_G f(y)g(y^{-1}x) dy$.
    \item \textbf{Natural Norm:} The $L^1$ norm, $\|f\|_1 = \int_G |f(x)| dx$.
    \item \textbf{Norm Closure (Banach *-algebra):} Completing $C_c(G)$ under the $L^1$ norm yields $L^1(G)$, the space of absolutely integrable functions. This represents the classical states (probability densities).
    \item \textbf{Weak Closure (von Neumann Algebra):} Representing $L^1(G)$ on $L^2(G)$ via left convolution and taking the double commutant yields $\mathcal{L}(G)$, the left group von Neumann algebra. This represents the full non-commutative algebra of quantum symmetries.
\end{itemize}

\paragraph{3. The Discrete Simplification}
When the group $G$ is a \textbf{discrete group} (such as the integers $\mathbb{Z}$ or the free group $F_2$), the continuous topological machinery dramatically simplifies into the algebra of sequences.
\begin{itemize}
    \item \textbf{The Measure:} The Haar integral reduces to a simple sum over group elements ($\int_G dx \to \sum_{x \in G}$).
    \item \textbf{The Seed Space:} The requirement of ``compact support'' simply means ``finite support.'' Therefore, the seed space $C_c(G)$ becomes exactly the algebraic group ring $\mathbb{C}[G]$, consisting of formal finite linear combinations of group elements.
    \item \textbf{The Classical Branch:} The space of continuous functions $C_0(G)$ becomes $c_0(G)$, the space of sequences that decay to zero. The von Neumann algebra $L^\infty(G)$ becomes $\ell^\infty(G)$, the space of all bounded sequences.
    \item \textbf{The Quantum Branch:} The Banach algebra $L^1(G)$ becomes $\ell^1(G)$, the space of absolutely summable sequences. Left convolution simplifies to a discrete, matrix-like product: $(f * g)(x) = \sum_{y \in G} f(y)g(y^{-1}x)$.
\end{itemize}

If the discrete group is further restricted to be a \textbf{finite group}, all notions of infinite sequences and topological closures vanish entirely. The limits collapse, and the spaces perfectly coincide as finite-dimensional vector spaces: $\mathbb{C}[G] \cong \ell^1(G) \cong \mathcal{L}(G)$.

The generalized Fourier transform serves as the mathematical bridge between these branches, intertwining the commutative geometry of $L^\infty(G)$ with the non-commutative dynamics of $\mathcal{L}(G)$ to form a complete Locally Compact Quantum Group.

\subsection[Example of D8]{Example of $D_8$: Non-Abelian Fourier Transform and Fusion Rules}

To make the algebraic structures concrete for a non-Abelian group, consider $D_8$, the dihedral group of symmetries of a square. The group has 8 elements, generated by a rotation $r$ (with $r^4=e$) and a reflection $s$ (with $s^2=e$ and $srs = r^{-1}$). The Hilbert space $L^2(D_8)$ is an 8-dimensional vector space spanned by the basis $\{\ket{g}\}_{g \in D_8}$.

As discussed previously, this vector space supports two distinct von Neumann algebras depending on the choice of multiplication:
\begin{enumerate}
    \item \textbf{The Classical Branch ($L^\infty(D_8)$):} If we equip the space with pointwise multiplication, $f \cdot g(x) = f(x)g(x)$, we obtain the commutative algebra $L^\infty(D_8)$. Operators in this algebra act diagonally in the group basis: $M_f \ket{g} = f(g)\ket{g}$.
    \item \textbf{The Quantum Branch ($\mathcal{L}(D_8)$):} If we equip the space with convolution, $f * g(x) = \sum_{y} f(y)g(y^{-1}x)$, we obtain the non-commutative group von Neumann algebra $\mathcal{L}(D_8)$. Because $D_8$ is a finite group, $\mathcal{L}(D_8) \cong L^1(D_8) \cong \mathbb{C}[D_8]$. This algebra is generated by the left regular representation $U_h \ket{g} = \ket{hg}$.
\end{enumerate}

Because $D_8$ is non-Abelian, the convolution algebra $\mathcal{L}(D_8)$ is non-commutative. The \textbf{Artin-Wedderburn theorem} dictates that any finite-dimensional semisimple algebra (like $\mathbb{C}[D_8]$) must decompose into a direct sum of matrix algebras. This decomposition is exactly the non-Abelian Fourier transform, which block-diagonalizes the convolution multiplication.

The group $D_8$ has five conjugacy classes: $\{e\}$, $\{r^2\}$, $\{r, r^3\}$, $\{s, r^2s\}$, and $\{rs, r^3s\}$. The number of irreducible representations equals the number of conjugacy classes. By the sum of squares formula $\sum d_i^2 = |G| = 8$, the only valid dimensions for the irreducible representations are four 1-dimensional representations (let us call them $1_A, 1_B, 1_C, 1_D$) and one 2-dimensional representation (denoted $2_E$). Thus, the Artin-Wedderburn decomposition gives:
\begin{align}
    \mathcal{L}(D_8) \cong \mathbb{C} \oplus \mathbb{C} \oplus \mathbb{C} \oplus \mathbb{C} \oplus M_2(\mathbb{C}).
\end{align}

\subsubsection*{Deriving Fusion Rules from the Center of $\mathcal{L}(D_8)$}

The fusion rules of the representation category $\mathrm{Rep}(D_8)$ govern how tensor products of representations decompose: $V_i \otimes V_j \cong \bigoplus_k N_{ij}^k V_k$. Interestingly, these fusion rules can be derived directly from the convolution algebra by looking at its center, $Z(\mathcal{L}(D_8))$.

The center consists of class functions (functions constant on conjugacy classes). Since there are five conjugacy classes, $Z(\mathcal{L}(D_8))$ is a 5-dimensional commutative subalgebra. A natural basis for the center is given by the normalized characters of the irreducible representations, $\chi_i(g) = \Tr(U^i(g))$.

By the orthogonality of characters, the characters form a basis of idempotents (up to normalization) under convolution. However, under \textit{pointwise multiplication}, the characters satisfy:
\begin{align}
    \chi_i(g) \cdot \chi_j(g) = \Tr(U^i(g)) \Tr(U^j(g)) = \Tr(U^i(g) \otimes U^j(g)).
\end{align}
Because the character of a tensor product representation is the pointwise product of the individual characters, decomposing $\chi_i \cdot \chi_j$ back into the irreducible character basis yields the fusion coefficients:
\begin{align}
    \chi_i \cdot \chi_j = \sum_k N_{ij}^k \chi_k.
\end{align}

Let us work out this detail for $D_8$. The four 1D representations act as signs ($\pm 1$) on the generators $r$ and $s$, and thus they form a $\mathbb{Z}_2 \times \mathbb{Z}_2$ subgroup under fusion. Their tensor products are trivial to compute: e.g., $1_B \otimes 1_C = 1_D$.
More interestingly, consider the fusion of the 2D representation $2_E$ with itself. Its character $\chi_E$ evaluates to $2$ on the identity $e$, $-2$ on $r^2$, and $0$ on all other elements. Pointwise squaring this character gives:
\begin{align}
    (\chi_E \cdot \chi_E)(e) = 4, \qquad (\chi_E \cdot \chi_E)(r^2) = 4, \qquad (\chi_E \cdot \chi_E)(\text{others}) = 0.
\end{align}
Notice that the sum of the characters of the four 1D representations ($1_A, 1_B, 1_C, 1_D$) yields exactly this function: each is $+1$ on $e$ and $+1$ on $r^2$, summing to $4$, while their signs perfectly cancel out on the remaining conjugacy classes. Therefore, we deduce the fusion rule:
\begin{align}
    2_E \otimes 2_E \cong 1_A \oplus 1_B \oplus 1_C \oplus 1_D.
\end{align}
This demonstrates how the commutative geometry of the center under pointwise multiplication encodes the non-commutative fusion dynamics of $\mathrm{Rep}(D_8)$.

\subsection{Parallel with operator algebraic quantum error correction}

The left/right regular decomposition also has a direct interpretation in
operator algebraic quantum error correction. In finite-dimensional
operator-algebra QEC, a correctable algebra has the canonical form
\begin{align}
    \mA_{\mathrm{log}}
    =\bigoplus_x B(\mH_{L,x})\otimes 1_{\mH_{G,x}}
    \subset
    B\left(\bigoplus_x \mH_{L,x}\otimes \mH_{G,x}\right),
\end{align}
where $\mH_{L,x}$ is the protected logical subsystem and $\mH_{G,x}$ is a
gauge or syndrome subsystem. The regular representation discussed above realizes precisely this structure with
\begin{align}
    \mH_{L,r}=\mH_r,\qquad
    \mH_{G,r}=\mH_r^*,\qquad
    \mA_{\mathrm{log}}=\mathcal L(G)
    =\bigoplus_{r\in\hat G}B(\mH_r)\otimes 1_{\mH_r^*}.
\end{align}
Note that in standard OAQEC, this structure typically describes a proper code subspace of the full physical Hilbert space. Here, the toy model treats the entire space $L^2(G)$ as the code space.
The commutant
\begin{align}
    \mA_{\mathrm{log}}'
    =\mathcal L(G)'
    =\bigoplus_{r\in\hat G}1_{\mH_r}\otimes B(\mH_r^*)
\end{align}
is the algebra of gauge or syndrome observables. Thus the left regular algebra
plays the role of the protected logical algebra, while the right regular algebra
plays the role of the complementary algebra.

The conditional expectation onto the left algebra is the Heisenberg-picture
version of erasing the gauge subsystem:
\begin{align}
    \mE_{\mathcal L}(X)
    =
    \bigoplus_{r\in\hat G}
    \left(\mathrm{id}_{B(\mH_r)}\otimes \tau_{\mH_r^*}\right)
    (P_r X P_r)\otimes 1_{\mH_r^*},
\end{align}
where $P_r$ projects onto $\mH_r\otimes\mH_r^*$ and
$\tau_{\mH_r^*}$ is the normalized trace on $B(\mH_r^*)$. This is the same map
as the average over the right regular action. Its predual sends a state to one
with the same reduced density matrix on the logical subsystem in each $r$-sector but with the multiplicity
factor replaced by the maximally mixed gauge state. In other words, the map
preserves expectation values of all logical operators in $\mathcal L(G)$ and
forgets the right-regular/gauge degrees of freedom.

This is the exact algebraic analogue of the information--disturbance viewpoint
on quantum error correction. For an ordinary channel, correctability of the
output is equivalent, up to quantitative constants in the approximate setting,
to the complementary channel being forgetful about the input. This fundamental decoupling duality between recoverability and forgetfulness of the environment was established in the OAQEC literature (e.g., B\'eny, Kempf, and Kribs (2007) and Kretschmann, Schlingemann, and Werner (2008)), and is discussed in modern contexts such as Section~2 of Hayden--Penington
\href{https://arxiv.org/pdf/1706.09434}{arXiv:1706.09434}. Here the ``environment'' is the commutant algebra $\mathcal L(G)'$: correcting
the logical algebra $\mathcal L(G)$ means that the complementary/right algebra
may contain the central sector label $r$, but it must not contain quantum
information inside the protected factor $\mH_r$.

The role of subspace error correction is also visible in the Wedderburn blocks.
Choosing a single irrep block $r$ gives a protected matrix algebra
$B(\mH_r)$, while $\mH_r^*$ is the corresponding gauge subsystem. Correcting
only a chosen block is analogous to decoding a chosen subspace. Correcting the
whole algebra $\mathcal L(G)$ means correcting all blocks at once, including the
central classical label $r$. Thus this setting gives a
group-theoretic toy model of the operator-algebraic QEC slogan:
\begin{align}
    \begin{gathered}
    \text{recoverable logical algebra}\\
    \Longleftrightarrow\\
    \text{Complementary algebra is forgetful of logical quantum information}.
    \end{gathered}
\end{align}


\section{Quantum dynamical systems: Group actions on von Neumann algebras}
\label{sec:group-action-vna-duality}

In the previous appendix, we studied groups through the lens of complex functions defined on them. In this subsection, we study groups through the lens of algebra-valued functions that live on them. Appendix~\ref{sec:group-case-appendix} discussed the regular representation and
the duality between the left and right group algebras. The corresponding
structure for a von Neumann algebra is obtained by replacing the trivial algebra
$\mathbb C$ with a von Neumann algebra $\mA$ equipped with an action of $G$.
Throughout this appendix, we take $G$ to be finite.



\subsection{Action, fixed points, and crossed product}

A group action on $\mA$ is a homomorphism
\begin{align}
    \alpha:G\to \mathrm{Aut}(\mA),\qquad
    \alpha_g\alpha_h=\alpha_{gh}.
\end{align}
With this structure, $\mA$ is a $G$-module algebra. Equivalently, it is a
module algebra over the Hopf algebra $\mathbb C[G]$, or in operator-algebra
language the triple $(\mA,G,\alpha)$ is a $W^*$-dynamical system. The point is
that the action is compatible with the algebra operations:
\begin{align}
    \alpha_g(ab)=\alpha_g(a)\alpha_g(b),
    \qquad
    \alpha_g(a^*)=\alpha_g(a)^*.
\end{align}
Equivalently, the group algebra acts on $\mA$ by
\begin{align}
    \rho_\alpha:\mathcal L(G)\otimes \mA\to \mA,\qquad
    \rho_\alpha(U_g\otimes a)=\alpha_g(a).
\end{align}
The invariant algebra is
\begin{align}
    \mA^G=\{a\in\mA:\alpha_g(a)=a,\ \forall g\in G\},
\end{align}
and the conditional expectation onto the invariant algebra is the group average
\begin{align}
    E_G(a)=\frac{1}{|G|}\sum_{g\in G}\alpha_g(a),\qquad
    E_G:\mA\to\mA^G.
\end{align}

When we have a group action $G$ on an algebra $\mathcal{A}$, in addition to $U_g$ (the regular representation of $G$ in $L^2(G)$), we have a new map that intertwines the action of the group on the algebra:
\begin{align}
    \lambda_g a=\alpha_g(a) \lambda_g
\end{align}
The crossed product $\mA\rtimes_\alpha G$ is the von Neumann algebra generated
by a copy of $\mA$ and unitaries $\{\lambda_g\}_{g\in G}$ implementing the
action:
\begin{align}
    \lambda_g a\lambda_g^*=\alpha_g(a),\qquad
    \lambda_g\lambda_h=\lambda_{gh}.
\end{align}
As a vector space, every element has a finite Fourier expansion
\begin{align}
    x=\sum_{g\in G}a_g\lambda_g,\qquad a_g\in\mA,
\end{align}
with twisted convolution product
\begin{align}
    (a_g\lambda_g)(b_h\lambda_h)
    =a_g\alpha_g(b_h)\lambda_{gh},
    \qquad
    (a_g\lambda_g)^*=\alpha_{g^{-1}}(a_g^*)\lambda_{g^{-1}}.
\end{align}
There is also a canonical conditional expectation back to $\mA$:
\begin{align}
    E_{\mA}\left(\sum_{g\in G}a_g\lambda_g\right)=a_e.
\end{align}
When $\mA=\mathbb C$, this crossed product reduces to the group algebra
$\mathcal L(G)$ of Appendix~\ref{sec:group-case-appendix}.

\subsection{Crossed Products as Twisted Left Convolution}

Let \(G\) be a discrete group. For a finitely supported scalar function
\(f:G\to \mathbb C\), the associated left convolution operator on
\(L^2(G)\) is
\[
  (L_f\xi)(k)
  =
  \sum_{g\in G} f(g)\xi(g^{-1}k).
\]
The composition of such operators is governed by convolution:
\[
  L_fL_h=L_{f*h},
  \qquad
  (f*h)(k)
  =
  \sum_{g\in G} f(g)h(g^{-1}k).
\]

Now let \(M\) be a von Neumann algebra and let
\(\alpha:G\to \operatorname{Aut}(M)\) be an action. The crossed product
\(M\rtimes_\alpha G\) is generated by \(M\) and unitaries \(u_g\) satisfying
\[
  \lambda_g a \lambda_g^*=\alpha_g(a).
\]
A finite element has the form
\[
  x=\sum_{g\in G} a_g \lambda_g,
  \qquad a_g\in M,
\]
so it may be regarded as an \(M\)-valued function on \(G\).

The product is computed by moving coefficients past group unitaries:
\[
  (a_g \lambda_g)(b_h \lambda_h)
  =
  a_g\alpha_g(b_h)\lambda_{gh}.
\]
Thus
\[
  (a *_\alpha b)(k)
  =
  \sum_{g\in G} a_g\,\alpha_g\!\left(b_{g^{-1}k}\right).
\]
This is precisely left convolution, except that the coefficients lie in
\(M\) and the group action twists the second factor.

\subsection{Dual coaction}

The crossed product has a canonical dual coaction. It is the map
\begin{align}
    \widehat\alpha:\mA\rtimes_\alpha G
    \to (\mA\rtimes_\alpha G)\,\bar\otimes\,\mathcal L(G)
\end{align}
defined on generators by
\begin{align}
    \widehat\alpha(a)=a\otimes 1,\qquad
    \widehat\alpha(\lambda_g)=\lambda_g\otimes U_g.
\end{align}
This is the direct analogue of the coaction
$\delta(U_g)=U_g\otimes U_g$ in Appendix~\ref{sec:group-case-appendix}. It keeps
track of the group label carried by each Fourier component $a_g\lambda_g$:
\begin{align}
    \widehat\alpha\left(\sum_{g\in G}a_g\lambda_g\right)
    =\sum_{g\in G}a_g\lambda_g\otimes U_g.
\end{align}
The fixed points of the dual coaction recover the original algebra:
\begin{align}
    (\mA\rtimes_\alpha G)^{\widehat G}
    :=\{x:\widehat\alpha(x)=x\otimes 1\}
    =\mA.
\end{align}
Thus the original action and the dual coaction exchange the two natural
operations:
\begin{align}
    \mA \longrightarrow \mA^G,
    \qquad
    \mA \longrightarrow \mA\rtimes_\alpha G.
\end{align}
The first map forgets charged operators by averaging over $G$, while the second
adjoins operators that implement the $G$-action. The expectations $E_G$ and
$E_{\mA}$ are the corresponding projections in these two pictures.

\subsection{Abelian and non-Abelian duals}

If $G$ is Abelian, the dual object $\hat G$ is again a group. The dual coaction
is equivalent to an ordinary action of $\hat G$ on the crossed product:
\begin{align}
    \widehat\alpha_\chi(a)=a,\qquad
    \widehat\alpha_\chi(\lambda_g)=\chi(g)\lambda_g,
    \qquad \chi\in\hat G.
\end{align}
Crossing again by this dual action gives the finite-group form of Takesaki
duality:
\begin{align}
    (\mA\rtimes_\alpha G)\rtimes_{\widehat\alpha}\hat G
    \simeq \mA\,\bar\otimes\,B(L^2(G)).
\end{align}
In other words, applying the symmetry extension and then the dual extension
returns the original algebra up to the universal matrix factor associated with
the regular representation.

For non-Abelian $G$, the set $\hat G$ of irreducible representations is not a
group, so there is no ordinary action of ``dual group elements'' on
$\mA\rtimes_\alpha G$. The correct dual object is the coaction above. After
Fourier transform, the $\mathcal L(G)$ leg decomposes into Wedderburn blocks
$B(\mH_r)$ labeled by $r\in\hat G$, and the dual sectors are matrix-valued
rather than one-dimensional characters. The tensor-product fusion rules of
$\mathrm{Rep}(G)$ enter through the coproduct
\begin{align}
    \Delta(U_g)=U_g\otimes U_g,
\end{align}
not through an action of a dual group.

\subsection{Dual inclusions}

The pair of inclusions
\begin{align}
    \mA^G\subset \mA,
    \qquad
    \mA\subset \mA\rtimes_\alpha G
\end{align}
are dual to one another. For outer actions on factors, both are finite-index
inclusions with index $|G|$. The inclusion $\mA^G\subset\mA$ remembers the
original symmetry through the charged sectors of $\mA$ as an $\mA^G$-bimodule,
whereas the inclusion $\mA\subset\mA\rtimes_\alpha G$ remembers the same data
through the implementing unitaries $\lambda_g$ and the dual coaction
$\widehat\alpha$. In the trivial case $\mA=\mathbb C$, this reduces to the
left/right regular duality of Appendix~\ref{sec:group-case-appendix}.

\subsection{Crossed product with $\hat{G}$}

In taking the crossed product $\mA\rtimes_\alpha G$, we introduced defect operators $\lambda_g$ that realize the action of the group on the algebra: $\alpha_g(a)=\lambda_g a \lambda_{g^{-1}}$. We take this action to be outer. This crossed product realizes the twisted convolution multiplication. An alternative choice of multiplication is point-by-point. In the regular representation of the group, we saw that this multiplication was realized using the projections $e_g=\ket{g}\bra{g}$: 
\begin{align}
    e_g e_h=\delta_{g,h}e_g, \qquad U_g e_h U_{g^{-1}}=e_{gh}\ .
\end{align}
To generalize point-by-point multiplication to algebra, we need a {\it co-action} map that pulls the generator of the action of $G$ on $\mA$ to that of the $L^2(G)$:
\begin{align}
    \hat{\alpha}(\lambda_g)=\lambda_g\otimes U_g\nn\\
    \hat{\alpha}(e_g)=1\otimes e_g\\
    \hat{\alpha}(a_g)=(a_g\otimes 1)\ .
\end{align}
Now, we can define the crossed product by $\hat{G}$: $(\mA\rtimes_\alpha G)\rtimes_{\hat{\alpha}}\hat{G}$ using the multiplication
\begin{align}
&X=\sum_{g,h}a_{g,h}\lambda_g e_h\nn\\
&\hat{\alpha}(X*Y)=\hat{\alpha}(X)\hat{\alpha}(Y)\nn\\
&(X*Y)_{k,s}=\sum_{gh=k}a_{g,hs}\alpha_g(b_{h,s})
\end{align}

\subsection{Physics interpretation}

We have the following inclusions:
\begin{align}
    \mA^G\subset \mA\subset \mA\rtimes_\alpha G\subset (\mA\rtimes_\alpha G)\rtimes_{\hat{\alpha}}\hat{G}\simeq \mA\otimes B(L^2(G))
\end{align}
The algebra $\mA$ is acted on by the symmetry transformation $\alpha$. The subalgebra $\mA^G$ is that of charge-neutral operators. The algebra $\mA$ has charged operators and charge-neutral operators. The conditional expectation (Haar integral over $\alpha$) projects out all charged operators. The crossed product $\mA\rtimes_\alpha G$ adds to $\mA$ (charged operators and charge neutral operators) defect operators $\lambda_g$. The dual symmetry $\hat{G}$ co-acts on $\mathcal{B}\equiv \mA\rtimes_\alpha G$. That is to say if we think of action $\alpha:G\times \mA\to \mA$, a co-action is $\hat{\alpha}:\mathcal{B}\to \mathcal{B}\times G$. The same way that $\mA^G$ is the fixed point of $\mA$ under the action of $G$, the algebra $\mA$ is the fixed point of $\mathcal{B}$ under the coaction $\hat{\alpha}$: that is $\hat{\alpha}(x)=x\otimes 1$ if and only if $x=a \lambda_{e}$. 

Gauging means to consider the invariant subalgebras of $\mA\rtimes_\alpha G$ under the action of the group. They are classified by $\sum_r a_r \lambda_r$, where $\lambda_r$ are the gauge-invariant defects, representing states of sector $r$. To describe the action of the dual symmetry on $\mathcal{B}$ we describe the action of the projection $e_r=\sum_g \ket{g}\bra{g}$ to the irreducible representation $r$
\begin{align}
    e_r\triangleright x=(\text{id}\otimes e_r)\hat{\alpha}(x)=x_r \lambda_r\otimes U_r=\hat{\alpha}(x_r\lambda_r)
\end{align}
So \(e_t\) projects onto the \(t\)-twisted sector. For a character \(\chi_r\) of an irreducible representation \(r\in\widehat G\),
\[
\chi_r(g)=\operatorname{Tr}_{H_r}\rho_r(g),
\]
we get
\[
\chi_r\triangleright (a_g\lambda_g)
=
\chi_r(g)\,a_g\lambda_g.
\]
Equivalently, using the normalized character \(\chi_r/d_r\),
\[
\frac{\chi_r}{d_r}\triangleright(a_g\lambda_g)
=
\frac{\chi_r(g)}{d_r}\,a_g\lambda_g.
\]
This is the operator-algebraic version of saying that a dual Wilson line labeled by \(r\) links a \(g\)-twisted defect and evaluates to the character \(\chi_r(g)\).

The relationship to the coaction is as follows. The coaction is the universal bookkeeping map
\[
a_g\lambda_g
\longmapsto
a_g\lambda_g\otimes U_g.
\]
The dual symmetry action is obtained by pairing the \(U_g\)-leg with functions or matrix coefficients:
\[
f(U_g)=f(g).
\]
Therefore
\[
(\operatorname{id}\otimes f)(a_g\lambda_g\otimes U_g)
=
f(g)a_g\lambda_g.
\]

If \(G\) is Abelian, the dual symmetry is an ordinary group \(\widehat G\). A character
\[
\chi\in\widehat G
\]
defines an automorphism
\[
\widehat\alpha_\chi(x)
=
(\operatorname{id}\otimes \chi)\delta(x).
\]
On generators,
\[
\widehat\alpha_\chi(a)=a,
\qquad
\widehat\alpha_\chi(\lambda_g)=\chi(g)\lambda_g.
\]

Thus, for
\[
x=\sum_g a_g\lambda_g,
\]
we have
\[
\widehat\alpha_\chi(x)
=
\sum_g \chi(g)a_g\lambda_g.
\]

This is an honest group action because for Abelian \(G\), characters multiply as
\[
(\chi\psi)(g)=\chi(g)\psi(g),
\]
and every irreducible sector is one-dimensional.

For non-Abelian \(G\), this is no longer an action by automorphisms. The irreducible characters \(\chi_r\) are not group-like elements. Instead, they satisfy
\[
\chi_r(g)\chi_s(g)
=
\chi_{r\otimes s}(g)
=
\sum_t N_{rs}^t\,\chi_t(g).
\]

Therefore the maps
\[
W_r(x):=(\operatorname{id}\otimes \chi_r)\delta(x)
\]
obey the fusion rule
\[
W_rW_s
=
\sum_t N_{rs}^t W_t.
\]

Thus the dual symmetry is not a group of invertible automorphisms. It is the non-invertible categorical symmetry
\[
\operatorname{Rep}(G).
\]

The simple dual symmetry operations are labeled by irreducible representations \(r\), and they fuse according to tensor-product fusion.

Now suppose the gauged algebra is the gauge-invariant twisted algebra
\[
B^G=(A\rtimes_\alpha G)^G,
\]
where the original group acts by conjugation:
\[
\beta_h(a_g\lambda_g)
=
\alpha_h(a_g)\lambda_{hgh^{-1}}.
\]

Elements of \(B^G\) have the form
\[
X_{[g]}(a)
=
\sum_{t\in G/C_G(g)}
\alpha_t(a)\lambda_{tgt^{-1}},
\qquad
a\in A^{C_G(g)}.
\]

These are gauge-invariant sums over the conjugacy class \([g]\).

A dual character \(\chi_r\) acts on this by
\[
W_r\bigl(X_{[g]}(a)\bigr)
=
\sum_{t\in G/C_G(g)}
\chi_r(tgt^{-1})\,\alpha_t(a)\lambda_{tgt^{-1}}.
\]

Since characters are class functions,
\[
\chi_r(tgt^{-1})=\chi_r(g),
\]
we get
\[
W_r\bigl(X_{[g]}(a)\bigr)
=
\chi_r(g)\,X_{[g]}(a).
\]

Thus, on gauge-invariant twisted sectors labeled by conjugacy classes \([g]\), the dual symmetry acts diagonally by irreducible characters:
\[
\boxed{
W_r\bigl(X_{[g]}(a)\bigr)
=
\chi_r(g)\,X_{[g]}(a).
}
\]

So the clean picture is
\[
\boxed{
\text{twisted sector }[g]
\quad\text{is measured by the dual symmetry }r
\quad\text{with eigenvalue }\chi_r(g).
}
\]

For Abelian \(G\), conjugacy classes are single elements and irreducible representations are characters. Then
\[
\widehat\alpha_\chi(\lambda_g)=\chi(g)\lambda_g
\]
is an ordinary invertible symmetry action.

For non-Abelian \(G\), the dual symmetry is the family of non-invertible operations
\[
W_r=(\operatorname{id}\otimes\chi_r)\delta,
\]
with fusion
\[
W_rW_s=\sum_t N_{rs}^tW_t.
\]

Thus the coaction
\[
\delta(\lambda_g)=\lambda_g\otimes U_g
\]
is the universal structure underlying the dual symmetry. The dual symmetry operations are obtained by probing the \(U_g\)-leg with functions, characters, or matrix coefficients.

\[
\boxed{
\text{dual invariants: }
(A\rtimes_\alpha G)^{\widehat G}=A
}
\]

\[
\boxed{
\text{dual crossed product: }
(A\rtimes_\alpha G)\rtimes_{\widehat\alpha}\widehat G
\cong
A\bar\otimes B(L^2(G))
}
\]

Thus ``gauging the dual symmetry gives back the original algebra'' means that the dual symmetry detects the \(\lambda_g\)-twisted sectors introduced by the first crossed product. Projecting to dual-invariant operators removes the nontrivial twisted sectors and leaves \(A\). Performing the full dual crossed product returns \(A\) up to the universal matrix factor \(B(L^2(G))\).

\section{Inclusions of algebras as generalization of groups}

Given an inclusion of 

\subsection{Inclusion of algebras of depth two}

Weak Hopf algebras to groups are Hopf algebras to groupoids.

An inclusion of factors with a trivial relative commutant is called {\it irreducible}. An irreducible inclusion of algebras with finite index is the same data as a Hopf algebra. 

\subsection{Inclusions of algebras of higher depth}

\section*{Appendix A: Algebraic and Measure-Theoretic Structures on Locally Compact Groups}
\label{app:haar_algebras}

In the construction of our operator algebraic models, it is necessary to formally distinguish the function spaces over the underlying symmetry group $G$. Let $G$ be a locally compact Hausdorff topological group equipped with a left-invariant Haar measure $\mu_L$. 

The translation of the left Haar measure by right multiplication defines the modular function $\Delta: G \to \mathbb{R}^+$, such that $d\mu_L(xg) = \Delta(g) d\mu_L(x)$. For the physical regimes considered in this work, the relevant symmetry groups (being compact, abelian, or semisimple Lie groups) are unimodular, meaning $\Delta(g) = 1$, and the measure is bi-invariant. However, keeping the modular function explicit clarifies the classical origins of the modular operator in Tomita-Takesaki theory, which later governs the emergence of dynamical time in the non-commutative von Neumann algebra setting.

The spaces $L^1(G)$ and $L^\infty(G)$ serve dual, yet algebraically distinct, roles representing states and observables, respectively:

\paragraph{The Observables ($L^\infty(G)$):}
The space of essentially bounded measurable functions, $L^\infty(G)$, forms a commutative $C^*$-algebra (and specifically, an abelian von Neumann algebra) under pointwise multiplication:
\begin{equation}
    (f \cdot g)(x) = f(x)g(x)
\end{equation}
Closure under this operation follows from the essential supremum, bounding the product as $\|fg\|_\infty \leq \|f\|_\infty \|g\|_\infty$. On the Hilbert space $L^2(G)$, elements of $L^\infty(G)$ act faithfully as diagonal multiplication operators.


\end{document}
